\acte{Les sept offrandes du Précieux Sang}
\usage{quot. ; surtout le ven. et en juillet, mois du Précieux Sang}
\rubrique{Après chaque offrande, on dit un Gloria Patri, puis l’oraison jaculatoire :}
\priere\textit{Béni soit à jamais et remercié Jésus, qui nous a sauvés par son Sang.}

\salut{I}{\textsc{Père Éternel,} je vous offre les mérites du Sang très précieux de Jésus, votre Fils bien-aimé et mon divin Rédempteur, pour la propagation et l’exaltation de ma mère chérie la sainte Église, pour la conservation et la prospérité de son chef visible, le Souverain Pontife romain, pour les cardinaux, les évêques, les pasteurs des âmes, et pour tous les ministres du sanctuaire.}
\salut{II}{Père Éternel, je vous offre les mérites du Sang très précieux de Jésus\,… pour la paix et la concorde entre les rois et les princes catholiques, pour l’humiliation des ennemis de la sainte foi, et pour la félicité du peuple chrétien.}
\salut{III}{Père Éternel, je vous offre les mérites du Sang très précieux de Jésus\,… pour le retour des incrédules, pour l’extirpation des hérésies, et pour la conversion des pécheurs.}
\salut{IV}{Père Éternel, je vous offre les mérites du Sang très précieux de Jésus\,… pour tous mes parents, mes amis et mes ennemis, les indigents, les infirmes, les affligés, et tous ceux pour qui vous savez que je dois prier et pour lesquels vous voulez que je prie.}
\salut{V}{Père Éternel, je vous offre les mérites du Sang très précieux de Jésus\,… pour tous ceux qui passeront aujourd’hui à l’autre vie, afin que vous les délivriez des peines de l’enfer, et que vous les admettiez très prochainement à la possession de votre gloire.}
\salut{VI}{Père Éternel, je vous offre les mérites du Sang très précieux de Jésus\,… pour ceux qui aiment un si grand trésor, pour ceux qui sont unis avec moi afin de l’adorer et de l’honorer, pour ceux enfin qui s’occupent à en propager la dévotion.}
\salut{VII}{Père Éternel, je vous offre les mérites du Sang très précieux de Jésus\,… pour tous mes besoins spirituels et temporels, pour le repos des saintes âmes du purgatoire, et spécialement de celles qui ont été plus dévotes à ce prix de notre rédemption et aux douleurs et aux peines de la très sainte Vierge, notre Mère chérie.}
\indulg{3 ans. — Plénière, aux conditions ordinaires, si les offrandes sont faites chaque jour pendant un mois. Pie VII, 22 sept. 1817 ; S.~Pénit., 12 mai 1931. — Ench.~1952, n\textsuperscript{o}~220. À la fin, on peut ajouter : \textit{Vive le Sang de Jésus, maintenant et toujours, et dans tous les siècles des siècles. Ainsi soit-il.}}
\culdelampe